\Askhsh{1742}{4}
\begin{ylh}{1}
    \item 3.4. 3ο Κριτήριο ισότητας τριγώνων
    \item 3.6. Κριτήρια ισότητας ορθογώνιων τριγώνων
    \item 3.7. Κύκλος - Μεσοκάθετος - Διχοτόμος
    \item 4.2. Τέμνουσα δύο ευθειών - Ευκλείδειο αίτημα
    \item 5.4. Ρόμβος
    \item 5.6. Εφαρμογές στα τρίγωνα
    \item 5.11. Ισοσκελές τραπέζιο.
\end{ylh}
Το τετράπλευρο $\mathrm{A}\mathrm{B}\Gamma\Delta$ του παρακάτω σχήματος είναι ρόμβος. Θεωρούμε $\mathrm{A}\mathrm{Z} \perp \Gamma\Delta$ και $\mathrm{A}\mathrm{E} \perp \Gamma\mathrm{B}$.
Να αποδείξετε ότι:
\begin{alist}
    \item Το τρίγωνο $\mathrm{Z}\mathrm{A}\mathrm{E}$ είναι ισοσκελές. \monades{6}
    \item Η ευθεία $\mathrm{A}\Gamma$ είναι μεσοκάθετος του τμήματος $\mathrm{Z}\mathrm{E}$. \monades{9}
    \item Αν $\mathrm{M}$ και $\mathrm{N}$ τα μέσα των πλευρών $\mathrm{A}\Delta$ και $\mathrm{A}\mathrm{B}$ αντίστοιχα, να αποδείξετε ότι $\mathrm{M}\mathrm{N} \parallel \mathrm{Z}\mathrm{E}$ και $\mathrm{Z}\mathrm{M} = \mathrm{E}\mathrm{N}$. \monades{10}
\end{alist}
% TODO: Προσθήκη σχήματος
